\section{Change of base field for algebraic preschemes}
\label{section:IV.4}

\oldpage[IV-2]{52}

\subsection{Dimension of algebraic preschemes}
\label{subsection:IV.4.1}

We shall develop in \S~5 the general theory of the dimension of preschemes.
The theory of the dimension of \emph{algebraic} preschemes, however, can be developed more elementarily and also has many special features; we therefore give here a rapid account independent of the general theory.

If $K$ is a field and $L$ an extension of $K$, we write $\operatorname{tr.deg}_K L$ for the \emph{transcendence degree} of $L$ over $K$.

\begin{definition}[4.1.1]
\phantomsection
\label{IV.4.1.1}
Let $k$ be a field and $X$ a prescheme locally of finite type over $k$.
The dimension of $X$ is the number
\begin{equation}
\dim(X)=\sup_x\operatorname{tr.deg}_k\boldsymbol{k}(x),
\tag{4.1.1.1}
\label{IV.4.1.1.1}
\end{equation}
where $x$ runs through the set of maximal points of $X$.
\end{definition}

We shall see in \S~5 \sref{IV.5.2.2} that Definition \sref{IV.4.1.1} depends only apparently on the base field $k$ over which $X$ is locally of finite type, and that the number $\dim(X)$ so defined agrees with the dimension of the \emph{underlying space} $X$ (defined in \sref[0]{0.14.1.2}).
Evidently,
\[
\dim(X)=\dim(X_{\mathrm{red}}).
\]

Each $\boldsymbol{k}(x)$ is an extension \emph{of finite type} of $k$ \sref[I]{I.6.3.3}, hence has finite transcendence degree over $k$.
If $X$ is of finite type over $k$ and nonempty, it is Noetherian \sref[I]{I.6.3.7}, hence has only finitely many irreducible components; consequently $\dim(X)$ is \emph{finite} and nonnegative.
Definition \sref{IV.4.1.1} gives
\[
\dim(\emptyset)=\setmin\infty.
\]

If $(X_\alpha)$ denotes the family of reduced closed sub-preschemes of $X$ whose underlying spaces are the irreducible components of $X$ \sref[I]{I.5.2.1}, then clearly
\begin{equation}
\dim(X)=\sup_\alpha\dim(X_\alpha).
\tag{4.1.1.2}
\label{IV.4.1.1.2}
\end{equation}
The calculation of the dimension is thus reduced to the case of \emph{integral} preschemes (locally of finite type over $k$).
Finally, we evidently have
\begin{equation}
\dim(X)=\dim(U)
\tag{4.1.1.3}
\label{IV.4.1.1.3}
\end{equation}
for every induced sub-prescheme on an \emph{everywhere dense} open subset $U$ of $X$.
The notion of dimension is thereby finally reduced to the case of a scheme \emph{affine of finite type over $k$}.

\begin{theorem}[4.1.2]
\phantomsection
\label{IV.4.1.2}
Let $X$ and $Y$ be two preschemes locally of finite type over a field $k$, and let $f:X\to Y$ be a $k$-morphism.

\oldpage[IV-2]{53}

\begin{enumerate}[(i)]
\item[\textup{(i)}] If $f$ is quasi-compact and dominant, then
\[
\dim(Y)\leq\dim(X).
\]
\item[\textup{(ii)}] If $f$ is quasi-finite, then
\[
\dim(X)\leq\dim(Y).
\]
\item[\textup{(iii)}] Suppose that $X$ is of finite type over $k$.
For $\dim(X)\geq n$ \textup{(respectively $\dim(X)\leq n$, respectively $\dim(X)=n$)}, it is necessary and sufficient that there exist an open subset $U$ dense in $X$ and a $k$-morphism
$g:U\longrightarrow\mathbf{V}(k^n)$
$(=\Spec(k[T_1,\ldots,T_n])$, which we shall also denote by $\mathbf{V}_k^n)$
that is surjective \textup{(respectively} finite, \textup{respectively} finite and surjective\textup{)}.
If $X$ is an affine scheme, one may take $U=X$.
\end{enumerate}
\end{theorem}

\begin{proof}
\emph{(i)} Let $y$ be a maximal point of $Y$.
The fibre $f^{-1}(y)$ contains a maximal point $x$ of $X$ \sref{IV.1.1.5}; hence $\boldsymbol{k}(x)$ is an extension of $\boldsymbol{k}(y)$, and therefore
\[
\operatorname{tr.deg}_k\boldsymbol{k}(y)
\leq
\operatorname{tr.deg}_k\boldsymbol{k}(x)
\leq
\dim(X).
\]
This proves (i).

\emph{(ii)} If $x$ is a maximal point of $X$, then $\boldsymbol{k}(x)$ is a finite extension of $\boldsymbol{k}(f(x))$ \sref[II]{II.6.2.2}, and hence has the same transcendence degree over $k$.
On considering the reduced closed sub-prescheme of $Y$ whose underlying space is $\overline{\{f(x)\}}$, we are reduced to the following corollary.
\end{proof}

\begin{corollary}[4.1.2.1]
\phantomsection
\label{IV.4.1.2.1}
Let $Y$ be a $k$-prescheme locally of finite type.
For every sub-prescheme $Z$ of $Y$, one has
\[
\dim(Z)\leq\dim(Y).
\]
Suppose in addition that all the irreducible components of $Y$ have the same dimension.
Then $Z$ is rare in $Y$ if and only if $\dim(Z)<\dim(Y)$.
\end{corollary}

\begin{proof}
Let $z$ be a maximal point of $Z$.
By considering one of the reduced closed sub-preschemes of $Y$ whose underlying space is an irreducible component of $Y$ containing $z$, we reduce to the case in which $Y$ is integral, with generic point $y$.
Then, by taking in $Y$ an affine open neighbourhood of $z$, we reduce to the case in which $Y$ is affine and of finite type over $k$, and $Z$ closed in $Y$.
Write
\[
Y=\Spec(A),
\qquad
Z=\Spec(A/\mathfrak a),
\]
where $A$ is an integral ring and $\mathfrak a$ is an ideal of $A$ distinct from $A$.

By the normalization lemma (Bourbaki, \emph{Alg.\ comm.}, chap.~V, \S~3, n\textsuperscript{o}~1, th.~1) there exists a finite sequence $(x_i)_{1\leq i\leq n}$ of elements of $A$, algebraically independent over $k$, such that $A$ is integral over $B=k[x_1,\ldots,x_n]$
and $\mathfrak a\cap B$ is generated by a subfamily $(x_i)_{1\leq i\leq p}$, possibly empty, of $(x_i)_{1\leq i\leq n}$.
Let $g:Y\longrightarrow\Spec(B)=\mathbf{V}_k^n$
be the finite dominant morphism corresponding to the canonical injection $B\to A$.
Since
\[
C=B/(\mathfrak a\cap B)
\cong k[x_{p+1},\ldots,x_n],
\]
$g$ induces on $Z$ a finite dominant morphism $Z\longrightarrow\mathbf{V}_k^{n-p}$.
Thus
\[
\operatorname{tr.deg}_k\boldsymbol{k}(y)=n,
\qquad
\operatorname{tr.deg}_k\boldsymbol{k}(z)=n-p,
\]
because $\boldsymbol{k}(y)$ and $\boldsymbol{k}(z)$ are finite extensions of $\boldsymbol{k}(g(y))$ and $\boldsymbol{k}(g(z))$, respectively.
This proves the first assertion.

If $p=0$, then necessarily $z=y$, since $y$ is the only point of $Y$ whose image under $g$ is the generic point of $\mathbf{V}_k^n$; in this case $Z$ contains a nonempty open subset of $Y$.
If on the contrary $p>0$, we necessarily have $z\neq y$, and $\overline{\{z\}}$ is therefore \emph{rare} in $\overline{\{y\}}$, which completes the proof of \sref{IV.4.1.2.1}.
\end{proof}

\begin{proof}[Proof of \sref{IV.4.1.2}, continued]
\emph{(iii)} The sufficiency of the stated conditions follows immediately from (i), (ii), and \sref{IV.1.5.4},~(v).
To prove necessity, choose a dense open subset $U$ of $X$ that is the union of pairwise disjoint irreducible affine open subsets $U_i$ $(1\leq i\leq m)$, each containing a maximal point of $X$ \sref[0\textsubscript{\mdseries I}]{0.2.1.6}.
We may further suppose that $X$ is reduced; if $X$ is affine, we may of course take $U=X$.

\oldpage[IV-2]{54}

The same argument as in \sref{IV.4.1.2.1} shows that, if $\dim(U_i)=n_i$, there exists a $k$-morphism $g_i:U_i\longrightarrow\mathbf{V}_k^{n_i}$
that is finite and dominant, hence surjective \sref[II]{II.6.1.10}.
Set $n=\dim(X)=\sup_i n_i$.
For every $n_i$ there is a morphism $h_i:\mathbf{V}_k^{n_i}\longrightarrow\mathbf{V}_k^n$
which is the closed immersion corresponding to the canonical homomorphism $k[T_1,\ldots,T_n]\longrightarrow k[T_1,\ldots,T_{n_i}]$
and is the identity when $n_i=n$.
Since $U$ is the disjoint union of the $U_i$, define $g$ on each $U_i$ to be $h_i\circ g_i$.
Then $g$ is finite and surjective.

When $n'\geq n$, composing $g$ with the canonical closed immersion $\mathbf{V}_k^n\longrightarrow\mathbf{V}_k^{n'}$
gives a finite morphism $U\to\mathbf{V}_k^{n'}$.
When $n'\leq n$, compose $g$ with the canonical morphism $p:\mathbf{V}_k^n\longrightarrow\mathbf{V}_k^{n'}$
corresponding to the canonical injection $k[T_1,\ldots,T_{n'}]\longrightarrow k[T_1,\ldots,T_n]$.
The morphism $p$ is faithfully flat, hence \emph{surjective}.
\end{proof}

\begin{remark}[4.1.3]
\phantomsection
\label{IV.4.1.3}
Corollary \sref{IV.4.1.2.1} shows that in formula \sref{IV.4.1.1.1} one may let $x$ run through the set of \emph{all points} of $X$.
\end{remark}

\begin{corollary}[4.1.4]
\phantomsection
\label{IV.4.1.4}
Let $X$ be a prescheme locally of finite type over a field $k$, and let $K$ be an extension of $k$.
Then $\dim(X\otimes_k K)=\dim(X)$.
\end{corollary}

\begin{proof}
It is enough to consider the case in which $X$ is of finite type over $k$.
The morphism $u:X\otimes_k K\longrightarrow X$
is faithfully flat \sref{IV.2.2.13},~(i).
Thus, if $U$ is an everywhere dense open subset of $X$, then $u^{-1}(U)=U\otimes_k K$
is dense in $X\otimes_k K$ \sref{IV.2.3.10}.
If $g:U\longrightarrow\mathbf{V}_k^n$
is finite and surjective, then so is $g_{(K)}:U\otimes_k K\longrightarrow\mathbf{V}_K^n$
\sref[I]{I.3.5.2} and \sref[II]{II.6.1.5}.
The corollary follows.
\end{proof}

\begin{corollary}[4.1.5]
\phantomsection
\label{IV.4.1.5}
Let $X$ and $Y$ be two preschemes locally of finite type over a field $k$.
Then $\dim(X\times_k Y)=\dim(X)+\dim(Y)$.
\end{corollary}

\begin{proof}
It is enough to prove the following.
Let $U$ and $V$ be affine neighbourhoods of points $x\in X$ and $y\in Y$ such that
\[
\operatorname{tr.deg}_k\boldsymbol{k}(x)=m=\dim(U),
\qquad
\operatorname{tr.deg}_k\boldsymbol{k}(y)=n=\dim(V).
\]
Then $\dim(U\times_k V)=m+n$.
Equivalently, by \sref{IV.4.1.2}, we may suppose that there exist finite surjective $k$-morphisms
\[
f:X\longrightarrow\mathbf{V}_k^m,
\qquad
g:Y\longrightarrow\mathbf{V}_k^n.
\]
The product morphism $f\times g:X\times_k Y\longrightarrow\mathbf{V}_k^{m+n}$
is finite and surjective \sref[I]{I.3.5.2} and \sref[II]{II.6.1.5}, and the result follows.
\end{proof}
